\section{Boundedness Analysis}\label{sec:boundedness}

\begin{definition}\label{def:bo}
  Given a method $C.m$, the \emph{boundness constraints for $C.m$}
  are the linear constraints over $\relevantpaths{C.m}$ (Def.~\ref{def:path}),
  using two versions for each program variable: $v$ is a variable in the pre-state
  and $v'$ is the same variable in the post-state.
  Such constraints are enriched with the $\max$ operator.
  The set of such constraints is written as $\bodenotations{C.m}$.
\end{definition}
%
Boundness constraits are an abstraction of denotations\ie they express how
the execution of a command modifies the boundness of some lists. Therefore, we provide
the boundness analysis as an abstract denotational semantics.
%
\begin{definition}[Definite Boundedness Analysis]\label{def:definite_boundness}
  Given a method $C.m$, the \emph{definite boundedness analysis} for commands
  is the map $\bo{C.m}{}:\Com\to\bodenotations{C.m}$ given
  in Fig.~\ref{fig:bo_commands}, where $C.m$ is dropped when clear
  from the context.
  %
  \begin{figure}[t]
    \begin{align*}
      \bo{}{\assign{v}{\new{\<List>}{}}}&=\{v'=0\}\cup U\\
      \bo{}{\assign{v_1}{\size{v_2}}}&=\{v_1'=v_2\}\cup U\\
      \underset{\text{(for any other $\exp$)}}{\bo{}{\assign{v}{\exp}}}&=\left\{v'=\pi\left|\begin{array}{l}
      %v\text{ has type \<List<$t$\texttt{>}> }\\
      \pair{\exp}{\dastates{}@\mathit{here}}\da{}\pair{\sigma'}{S}\\
      \pi\in S\cap\relevantpaths{}
      \end{array}\right.\right\}\cup U\\
      \bo{}{\assign{\field{v_1}{f}}{v_2}}&=\left\{(\pi_1.f)'=\pi_2\left|\begin{array}{l}
      \pi_1\in\dastates{}@\mathit{here}(v_1)\\
      \pi_2\in\dastates{}@\mathit{here}(v_2)\cap\relevantpaths{}\\
      \pi_1.f\in\relevantpaths{}
      \end{array}\right.\right\}\cup U\\
      \bo{}{\readmap{v}{v}{v}}&=U\\
      \bo{}{\require{v_1}{v_2}}&=\{v_1\ \mathit{cop}\ v_2\}\cup U\\
      \bo{}{\push{v_1}{v_2}}&=
      \left\{\pi'=\pi+1\left|\begin{array}{l}
      \pi\in\dastates{}@\mathit{here}(v_1)\\
      \text{\<[any]> does not occur in $\pi$}
      \end{array}\right.\right\}\\
      &\qquad\cup\left\{\pi'=\max(\pi,v_1+1)
      \left|\begin{array}{l}
      \pi\in\dastates{}@\mathit{here}(v_1)\\
      \text{\<[any]> occurs in $\pi$}\\
      \pi\survives{\push{v_1}{v_2}}
      \end{array}\right.\right\}\\
      &\qquad\cup\left\{\pi'=\pi\in U\left|\begin{array}{l}
      \tau(v_1)=\text{\<List<$t'$\texttt{>}>}\\
      \text{the type of $\pi$ is \<List<$t''$\texttt{>}>}\\
      \text{and $t'$ and $t''$ have no shared subtype}
      \end{array}\right.\right\}\\
      \bo{}{\receive{v}{v}}&=U\\
      \bo{}{\call{v}{m}{v^*}}&=U\\
      \bo{}{\ifetwo{v_1}{v_2}{\com_1}{\com_2}}&=\left(\{v_1\ \mathit{cop}\ v_2\}
      \cup\bo{}{\com_1}\right)\oplus\bo{}{\com_2}\\
      \bo{}{\for{v_1}{v_2}{\com}}&=\lfp.F\mapsto\mathit{Id}\oplus(\bo{}{\com}\otimes F)\\
      \bo{}{\sequence{\com_1}{\com_2}}&=\bo{}{\com_1}\otimes\bo{}{\com_2}
    \end{align*}
    \caption{The definite boundness analysis of the commands. It uses the definite
      aliasing analysis at the command on the left-hand side of each equation,
      written as $\dastates{}@\mathit{here}$. The \emph{frame condition} $U$ (for \emph{unchanged})
      states that no existing relevant paths is affected, unless it is modified by side-effect.
      For instance, for the first equation,
      it is $U=\{\pi'=\pi\mid\pi\in\relevantpaths{}\text{ and }\pi\survives{\assign{v}{\new{\<List>}{}}}\}$.}
    \label{fig:bo_commands}
  \end{figure}
\end{definition}

\textbf{This is a sort of simplified path-length. The simplifying
  assumptions must be clarified. The idea is that one does not need
  the supporting analyses such as non-cyclicity or sharing, since the
  domain of analysis is very specific and lists are very well encapsulated.}

The static analysis infers size change constraints about the containers
reachable from the state of the smart contracts, for each public API
method or constructor of the smart contract.

For the example in Fig.~\ref{fig:ponzi1}, the inferred approximations are:
%
\begin{description}
\item[\<Ponzi1()>:] $\<investors>'=1$
\item[\<invest()>:] $\<investors>'=\<investors>+1$
\end{description}
%
For the esample in Fig.~\ref{fig:ponzi2} they are:
%
\begin{description}
\item[\<Ponzi2()>:] $\<investors>'=1$
\item[\<invest()>:] $\<investors>'=\<investors>+1$
\item[\<withdraw()>:] $\<investors>'=\<investors>$
\end{description}
%
For the example in Fig.~\ref{fig:ponzi3} they are:
%
\begin{description}
\item[\<Ponzi3()>:] $\<investors>'=1$
\item[\<invest()>:] $\<investors><100\wedge\<investors>'=\<investors>+1$
\end{description}
%
For the example in Fig.~\ref{fig:lottery1} they are:
%
\begin{description}
\item[\<Lottery1()>:] $\<rounds>'=1\wedge\<rounds[any].sales>'=0$
\item[\<buy()>:] $\<rounds>\le\<rounds>'\le\<rounds>+1\wedge\<rounds[any].sales>\le\<rounds[any].sales>'\le\<rounds[any].sales+1>$
\item[\<drawWinner()>:] $\<rounds>'=\<rounds>\wedge\<rounds[any].sales>'=\<rounds[any].sales>$
\end{description}
%
For the example in Fig.~\ref{fig:lottery2} they are:
%
\begin{description}
\item[\<Lottery2()>:] $\<rounds>'=1\wedge\<rounds[any].sales>'=0$
\item[\<buy()>:] $\<rounds>\le\<rounds>'\le\<rounds>+1\wedge\<rounds[any].sales>\le\<rounds[any].sales>'\le\max(5000,\<rounds[any].sales>)$
\item[\<drawWinner()>:] $\<rounds>'=\<rounds>\wedge\<rounds[any].sales>'=\<rounds[any].sales>$
\end{description}
